Based on the modal analysis we ran simulations for the frequency range \SIrange{1}{900}{Hz} and tracked the maximum displacement, as is shown in figure \ref{fig:max-displacement-6}.
At low frequencies, up to about \SI{400}{Hz}, we observed a linear decline of about 2 decade/decade, or -\SI{60}{\deci\bel\per decade} across the entire device.
As frequency increased to \SI{850}{Hz} we saw the non-linear behaviour, such that the displacement of each sensor relative to one another increased.
Despite the non-clinical frequencies in which it existed, this change in relative movement could be a tool, therefore they were explored in further detail.

\begin{figure}[h]\centering
    \begin{tikzpicture}
        \begin{loglogaxis}[
            % width = 14cm, height = 6cm,
            xlabel = {Frequency [Hz]}, ylabel = {Maximum displacement [m]},
            title = {},
        ]
            \addplot+[only marks] table [ col sep=comma,]
                {Modelling/Frequencies/spigot_maxima.csv};
            \addlegendentry{Spigot}
            \addplot+[ only marks ] table [col sep=comma,]
                {Modelling/Frequencies/stem_maxima.csv};
            \addlegendentry{Stem}
            \addplot+[ only marks ] table [col sep=comma, ]
                {Modelling/Frequencies/collar_maxima.csv};
            \addlegendentry{Collar}
        \end{loglogaxis}
    \end{tikzpicture}

\caption{Maximum displacement as a function of frequency for 6 periods.}
\label{fig:max-displacement-6}
\end{figure}

\paragraph{Considerations}
The data points for low frequencies were sparse, and could have been 10 or 20 divisions per decade.
This is only possible because the creation of the simulations and the data processing were completely automated with Python scripts.
Ultimately this was not an issue because of the linear behaviour observed, but it should have been a design consideration.

\subsection{Relative displacements}
Figure \ref{fig:normalisation} characterises the maximum displacement for all parts relative to the maximum displacement of the spigot.
For the clinically relevant frequencies, the displacement of the stem hovers nicely around 4 times that of the spigot. 

Whether the large relative displacement observed around \SI{700}{Hz} is useful is yet to be ascertained.
We could take advantage of the large relative displacement if it translates to larger forces in subsequent studies, though it likely comes with a loss in resolution \cite{rosenbaum-chou_development_2016}.
For this reason we continued to characterise vibrations beyond what is clinically relevant.
% \clearpage
\begin{figure}[h] \centering
\begin{tikzpicture}[scale=0.95]
    \begin{axis}[
        % width = 14cm, height = 6cm,
        xlabel = {Frequency [Hz]}, ylabel = {Displacement relative to spigot},
        % legend pos = north west,
        ymin=-1,
        extra y ticks = {1},
    ]
        \addplot+ [only marks] table [col sep=comma]
            {Modelling/Frequencies/spigot_maxima_normalised.csv};
        \addlegendentry{Spigot}
        \addplot+[only marks, ] table [col sep=comma,]
            {Modelling/Frequencies/stem_maxima_normalised.csv};
        \addlegendentry{Stem}
        \addplot+[ only marks,] table [col sep=comma,]
            {Modelling/Frequencies/collar_maxima_normalised.csv};
        \addlegendentry{Collar}
    \end{axis}
\end{tikzpicture}
\caption{Maximum displacement relative to spigot as a function of frequency.}
\label{fig:normalisation}
\end{figure}
% \clearpage